\documentclass{article}
\usepackage{xunicode, amsmath, amssymb}
\usepackage[utf8]{inputenc}
\usepackage[none]{hyphenat}

\title{Every Function is a Polynomial: the Proof}
\author{Shaunak Desai}
\date{16 August 2025}

\begin{document}

\maketitle
\section{Recap}

Last week, we covered a heuristic formulation of Taylor series, where we may expand any suitably nice function $f(x)$ around a point $x=a$ via an infinite series of polynomials:

$$
f(x) = \sum_{n=0}^{\infty} \frac{f^{(n)}(a)}{n!} (x-a)^n
$$

The Taylor series around $x=0$ is often known as the Maclaurin series, after the Scottish mathematician Colin Maclaurin, who used them in his \textit{Treatise of Fluxions} to analyse infinitely differentiable functions.

(Incidentally, Colin Maclaurin holds the record for the youngest maths professor in history, gaining the post of Professor of Mathematics at the University of Aberdeen at the age of 19. He held the record for youngest professor for almost three centuries, until it was broken in 2008 by Alia Sabur at the age of 18.)

Some quick, but important, addendums:
\begin{enumerate}
    \item For a function's Taylor series to exist at a point $x=a$, it must be infinitely differentiable at that point. This may be obviously evident through the definition, but it is worth mentioning.
    \item The so-called "0th derivative" of any function is taken to be the function itself; think of the 0th derivative as not taking the derivative in the first place. When the Taylor series is written in sigma notation, this corresponds to the $n=0$ term being the function itself evaluated at the point of expansion, i.e. $f(a)$.
    \item Perhaps counterintuitively, a function's Taylor series need not necessarily converge to the function itself. A function $f(x)$ is \textbf{analytic} at a point $x=a$ if its Taylor series around $x=a$ converges to $f(x)$. For example, the function
    $$
    f(x) =
    \begin{cases}
    e^{-1/x^2} & \text{if } x \neq 0 \\
    0 & \text{if } x = 0
    \end{cases}
    $$
    is infinitely differentiable, but not analytic, at $x=0$. Evaluate some derivatives at $x=0$ to see why!
\end{enumerate}

\end{document}
